%------------------------------------------------------------------------------%
\begin{exercises}

%------------------------------------------------------------------------------%
\begin{exercise}
Calcule as integrais, usando integração por partes
\begin{mathlist}
  \item \int x \cos(5x)dx
  \item \int xe^{-x}dx
  \item \int t^2 \sin(\beta t) dt\) onde $\beta$ é uma constante.
  \item \int (x^2+2x) \cos(x) dx
  \item \int p^5 ln(p) dp
  \item \int (ln(x))^2 dx
  \item \int  t \sec^2(2t) dt
  \item \int e^{2x} \sin(3x) dx
  \item \int_4^9 \frac{\ln(y)}{\sqrt{y}}
  \item \int_0^{\frac{1}{2}} (x \cos(\pi x))dx
  \item \int_0^1 (x^2+1)e^{-x}dx
  \item \int_1^2 x^4 (ln(x))^2 dx
  \item \int_1^{\sqrt{3}} \arctan(\frac{1}{x})dx
  \item \int x \tan^2(x) dx$. Dica: $\tan^2(x)=1-\sec^2(x)
\end{mathlist}
\begin{answer}
  \begin{mathlist}
    \item \frac{1}{5}x \sin(5x)+\frac{1}{25} \cos(5x)+c
    \item -xe^{-x}-e^{-x}+C
    \item -\frac{1}{\beta}t^2 \cos(\beta t)+\frac{2}{\beta^2}t \sin(\beta t)+ \frac{2}{\beta ^3}\cos(\beta t)+c
    \item (x^2+2x) \sin(x)+(2x+2)\sin(x)-2 \sin(x)+C
    \item \frac{1}{6}p^6 ln(p) +\frac{1}{36}p^6+C
    \item x(ln(x))^2-2x ln(x)+2x+c
    \item \frac{1}{2} t \tan(2t) - \frac{1}{4} ln( \vert \sec(x) \vert)+C
    \item \frac{1}{13}e^{2x} \left( 2 \sin(3x)-3 \cos(3x)+C\right)
    \item  6ln(9)-4ln(4)-4
    \item \frac{\pi-2}{\pi^2}
    \item -\frac{6}{e}+3
    \item \frac{32}{5}(ln2)^2- \frac{64}{25}ln2+ \frac{62}{126}
    \item \frac{\pi \sqrt{3}}{6}- \frac{\pi}{4}+ \frac{1}{2}ln2
    \item x \tan(x)- ln(\vert \sec(x) \vert)- \frac{1}{2}x^2+C
  \end{mathlist}
\end{answer}
\end{exercise}

%------------------------------------------------------------------------------%
\begin{exercise}
Encontre a área sob a curva $y=x \sin(x)$ de $0$ a $\pi$.
\begin{answer}
  $A=\pi$
\end{answer}
\end{exercise}

%------------------------------------------------------------------------------%
\begin{exercise}
Uma partícula se move ao longo de uma reta com velocidade igual a
$v(t)=t^2e^{-t}$ metros por segundo após $t$ segundos.
Qual a distância que essa partícula percorrerá nos primeiros $t$ segundos?
\begin{answer}
  $s=2-e^{-t}(t^2+2t+2)m$
\end{answer}
\end{exercise}

%------------------------------------------------------------------------------%
\begin{exercise}
Calcule as integrais
\begin{mathlist}
  \item \int x \arctan(x) dx
  \item \int_{1}^{2} \frac{(\ln(x))^2}{x^3}dx
  \item \int x^3 e^x dx
  \item \int e^{-y} \cos(2y)dy
  \item \int_0^1 \frac{y}{e^{2y}}dy
\end{mathlist}
\end{exercise}

%------------------------------------------------------------------------------%
%\begin{exercise}
%\begin{answer}
%\end{answer}
%\end{exercise}

%------------------------------------------------------------------------------%
%\begin{exercise}
%\begin{mathlist}[2]
%  \item
%  \item
%  \item
%\end{mathlist}
%\begin{answer}
%  \begin{alphalist}[3]
  %    \item
  %    \item
  %    \item
  %  \end{alphalist}
%\end{answer}
%\end{exercise}

\end{exercises}
%------------------------------------------------------------------------------%
